\chapter{Background and Related Work}
\label{ch:background}

\section{Fall detection and application context}
Falls are an important source of injury and loss of independence among older adults. A wearable detector is intended to identify a suspected event when a person is away from fixed cameras or room sensors, preserve a traceable event record, and support a caregiver's follow-up. Such a detector is an assistive alerting tool; it does not diagnose injury, replace emergency services, or establish that a person is safe when no alarm occurs. WHO guidance motivates the broad fall-prevention context, and Nepal's census-based population projections provide demographic planning context; neither source measures local fall prevalence or demand for SPARK \cite{who2007falls,cbs2021nepal}.

The research literature spans threshold rules, feature-based classifiers, temporal neural networks, edge inference, and systems that send candidate signals to a server. Each choice affects more than classification accuracy: sensor placement, event definitions, false alarms over ordinary activity, compute and radio duty, notification delay, usability, and evaluation population all matter. SisFall is a public benchmark for reproducible development, but its waist-mounted signals and largely acted falls do not validate a dorsal-wrist system or natural falls \cite{sucerquia2017sisfall}.

\section{Comparative literature review}
Threshold methods are computationally simple and interpretable. Kangas et al. compared low-complexity algorithms at waist, wrist, and head placements; their results emphasize that the best site and rule depend on placement, and that the wrist signal poses a particular challenge \cite{kangas2008comparison}. Bourke et al. evaluated tri-axial thresholds using simulated falls and older adults' daily activities, illustrating both the value of trunk signals and the limits of transferring a threshold across sensor positions \cite{bourke2008threshold}. These studies motivate a gate as a candidate filter, not the assumption that a fixed gate preserves every fall.

Learned models can represent temporal patterns that a single magnitude rule cannot. SisFall provides 15 fall activities and 19 activities of daily living from 38 participants, with inertial sensors; it is a useful public source for offline model development, subject-grouped evaluation, and comparison of preprocessing choices \cite{sucerquia2017sisfall}. CNN and recurrent approaches represent temporal structure differently. Nho et al.'s cited title and the listed Wang et al. citation could not be matched to authoritative primary records using the supplied metadata; this report therefore does not repeat their performance claims or use them to rank architectures. Those entries remain visible in the proposal-reference audit.

Embedded inference literature establishes a feasible runtime route, not a guarantee for this board and model. TensorFlow Lite Micro describes an interpreter and evaluation for constrained embedded systems; actual RAM, timing, operator coverage, and whole-firmware size still require target measurements \cite{david2021tensorflow}. Xu et al. provide a directly relevant hybrid example: their wearable threshold stage forwards candidate acceleration windows to a server-side CNN, so the classifier is not fully local to the wearable \cite{xu2021fusion}. Yu et al. report a TinyCNN wearable prototype and an Arduino Nano BLE 33 Sense implementation. Its published record is in Biomedical Signal Processing and Control, volume 86, article 105325 (2023), not the Measurement citation given in the proposal \cite{yu2023tinycnn}. Hu et al. report MicroFallNet on UMAFall and FallAllD, including an ESP32 implementation; the correct record is Pervasive and Mobile Computing, volume 109, article 102046 (2025), not volume 107/article 101980 \cite{hu2025microfallnet}.

Other system designs broaden the comparison beyond classifier topology. Tseng et al. describe a finite-state wearable detector with GPS, NB-IoT, and notifications, demonstrating a location-aware communication path distinct from SPARK's local gateway design \cite{sensors2025wearable}. Bhattacharyya et al. describe a CNN-kNN fall detector on a RISC-V system implemented on FPGA; this is an accelerator-oriented hardware point of comparison, not evidence that a low-cost MCU has equivalent resources \cite{vigil2024fpga}. Ba\~nos et al.'s mHealthDroid is relevant to mobile-health framework context, but SPARK's phone client is deliberately display-only and does not adopt its broader sensor integration architecture \cite{banos2014mhealthdroid}.

Explainability terminology also constrains the comparison. SHAP defines additive feature attributions grounded in Shapley values \cite{lundberg2017shap}. The SPARK gateway's input-times-gradient calculation is a different attribution method, while the live v1 event's peak-feature display is a heuristic summary because the packet does not contain a waveform. Neither should be labeled SHAP. Yao et al. study ECG heartbeat interpretation using CNN and GRU; that biomedical time-series example supplies context but is not fall-detection evidence \cite{yao2021interpretation}. Liu et al.'s IEEE conference paper concerns explainable wearable human-activity recognition rather than fall-event explanation \cite{liu2024explainable}. Kundu et al.'s review is an arXiv preprint submitted in 2026 (arXiv:2604.09799), rather than a 2025 peer-reviewed article as the proposal bibliography says \cite{kundu2025explainable}.

\begin{table}[htbp]
\centering\small
\caption{Comparison of verified and directly relevant approaches}
\label{tab:literature-comparison}
\begin{tabular}{p{2.4cm}p{3.0cm}p{3.3cm}p{3.4cm}}
\hline
Approach & Evidence in cited work & Design implication & SPARK status and boundary \\
\hline
Threshold rules & Low-complexity inertial algorithms; placement matters \cite{kangas2008comparison,bourke2008threshold} & Measure gate sensitivity and false alarms at the intended placement & Gate code exists; its 2.5 g / greater-than-300 ms behavior awaits event-level evaluation \\
Public inertial benchmark & SisFall contains falls and daily activities from 38 people \cite{sucerquia2017sisfall} & Use grouped splits and disclose sensor placement and acted-fall limits & Offline 100 Hz, 200-sample, six-channel windows; wrist transfer remains unvalidated \\
Embedded runtime & TFLite Micro targets constrained embedded inference \cite{david2021tensorflow} & Measure linked target memory, operators, and latency & Local model bytes are known; MCU inference remains a stub \\
Hybrid gate plus CNN & Xu et al. send candidate windows to a server CNN \cite{xu2021fusion} & Report compute location and communication dependency explicitly & SPARK intends local inference, but has not demonstrated it on hardware \\
Wearable edge CNN & TinyCNN and MicroFallNet report microcontroller implementations \cite{yu2023tinycnn,hu2025microfallnet} & Treat published results as feasibility context, not a head-to-head comparison & Different datasets, placements, models, and protocols prevent direct score comparison \\
Location and FPGA systems & Tseng et al. include localization/notification; Vigil uses FPGA-oriented RISC-V hardware \cite{sensors2025wearable,vigil2024fpga} & Compare system goals and hardware classes, not only accuracy & Context only; neither architecture is the SPARK implementation \\
Explanation methods & SHAP has Shapley-value semantics; HAR/ECG work covers different tasks \cite{lundberg2017shap,liu2024explainable,yao2021interpretation,kundu2025explainable} & Name the actual attribution method and input used & Gradient-times-input and peak-feature fallback are distinct; explanation validity is pending \\
\hline
\end{tabular}
\end{table}

\section{Hardware and software theory}
An inertial measurement unit combines accelerometer and angular-rate measurements. A three-axis accelerometer observes specific force in its sensor frame; a gyroscope observes angular velocity. Orientation, gravity separation, scale factors, axis alignment, and units must be considered before thresholds or models trained on one representation are applied to another. The current preprocessing selection retains raw SisFall sensor counts, whereas the event protocol describes peak values in g and degrees per second. The exact transform linking those representations is unresolved and must be specified before firmware integration.

The intended embedded design uses an ESP32-S3 and MPU6050. Sensor acquisition, timestamping, buffering, preprocessing, gate evaluation, model invocation, and event delivery form a real-time chain. A gateway receives and stores events, generates local reports, and offers a display API. The phone is a display client. The architecture is staged to make failures inspectable; the current sensor reads, TFLite invocation, transport writes, and logger loop are not complete physical implementations.

For six input channels, a temporal convolution computes learned local filters over adjacent samples. Pooling reduces temporal resolution, global average pooling summarizes activations, and a dense output layer produces two class scores. INT8 quantization represents a real value approximately as $x=s(q-z)$ using model-specific scale $s$ and zero point $z$. This reduces weight storage but introduces quantization error and does not account for runtime buffers or the rest of the firmware. These are engineering mechanisms, not performance guarantees.

\section{Interpretation and study position}
SPARK's defensible present contribution is an inspectable software pipeline, offline model artifacts, and an evaluation design that separates archived results from future integrated measurements. The literature supports testing threshold filtering, quantized embedded inference, event traceability, and local explanation as design questions. It does not establish novelty, clinical suitability, power savings, or superiority for this system. The exact source records and claim-level disposition for all twenty proposal references are documented in \texttt{docs/thesis\_work/PROPOSAL\_COVERAGE\_AND\_REFERENCE\_AUDIT.md}.
